% =====================================================================
% Thesis — Chapter 2
% Statistical Preliminaries
% Class: apa7 (doc option) + biblatex-apa (APA 7th)
% Compile: pdflatex -> biber -> pdflatex -> pdflatex
% =====================================================================
\documentclass[doc,11pt,donotrepeattitle]{apa7}

\usepackage[american]{babel}
\usepackage{csquotes}
\usepackage{amsmath}
\usepackage{amssymb}
\usepackage{booktabs}
\usepackage{array}
\usepackage{ragged2e}

\usepackage[style=apa,sortcites=true,sorting=nyt,backend=biber]{biblatex}
\DeclareLanguageMapping{american}{american-apa}
\addbibresource{bibliografia.bib}

\newcolumntype{L}[1]{>{\RaggedRight\arraybackslash}p{#1}}

% apa7 loads hyperref; hide the colored link borders for a clean print look
\hypersetup{hidelinks}

\title{Statistical Preliminaries}
\shorttitle{Statistical Preliminaries}
\author{Vittorio Guerrieri}
\affiliation{Chapter~2 of the thesis}

\begin{document}

\maketitle

% =====================================================================
\section{Orientation}
% =====================================================================

The methods developed in the following chapters draw on a handful of
statistical ideas that, while standard in statistics and machine learning, are
not all part of the routine toolkit of a working psychologist. This chapter
collects them in one place and explains them at the level of detail needed to
follow the rest of the thesis. It is deliberately self-contained and tutorial in
tone: a reader comfortable with these concepts may skim it, while a reader who
last met some of them in an undergraduate methods course should find here
enough to read Chapters~3 to 5 without consulting outside references.

The selection is driven by what the method actually uses. The
correlation coefficient (Section~\ref{sec:corr}) is the raw material of the
ReReReRe detector, which compares a respondent's \emph{individual-level}
coherence with a baseline. Exploratory factor analysis
(Section~\ref{sec:efa}) supplies the latent-variable model that justifies why
attentive responses to different items covary in the first place, and parallel
analysis---the procedure ReReReRe uses to estimate the number of factors---is a
factor-retention method. Permutation tests (Section~\ref{sec:perm}) are the
inferential engine behind ReReReRe's per-respondent z-score, and the bootstrap
(Section~\ref{sec:boot}) provides the confidence intervals used to quantify the
detector's contribution. The Random Forest (Section~\ref{sec:rf}) is the
classifier that combines ReReReRe with the auxiliary detectors into an
ensemble, cross-validation (Section~\ref{sec:cv}) is how its performance is
estimated honestly, and the battery of evaluation metrics
(Section~\ref{sec:metrics})---above all the Matthews correlation
coefficient---is how detection quality is reported throughout.

% =====================================================================
\section{Correlation and Covariance}
\label{sec:corr}
% =====================================================================

Two variables are \emph{correlated} when knowing one tells us something about
the other. The workhorse measure of linear association is the Pearson
product--moment correlation coefficient \parencite{pearson1896mathematical}.
For two variables $X$ and $Y$ measured on $n$ units, it is defined as the
covariance of $X$ and $Y$ divided by the product of their standard deviations,
\begin{equation}
r_{XY} \;=\; \frac{\operatorname{cov}(X,Y)}{s_X\, s_Y}
\;=\; \frac{\sum_{i=1}^{n}(x_i-\bar{x})(y_i-\bar{y})}
{\sqrt{\sum_{i=1}^{n}(x_i-\bar{x})^2}\,\sqrt{\sum_{i=1}^{n}(y_i-\bar{y})^2}} .
\label{eq:pearson}
\end{equation}
The coefficient ranges from $-1$ (perfect negative linear association) through
$0$ (no linear association) to $+1$ (perfect positive linear association). The
numerator, the \emph{covariance}, is positive when above-average values of $X$
tend to co-occur with above-average values of $Y$; dividing by the standard
deviations rescales this to the unit interval, making correlations comparable
across variables measured on different scales.

Three properties of Equation~\ref{eq:pearson} matter for what follows. First,
$r$ captures only \emph{linear} association: two variables can be strongly
related in a curved pattern yet have $r \approx 0$. Second, $r$ is invariant to
separate linear rescalings of $X$ and $Y$: adding a constant to every value, or
multiplying every value by a positive constant, leaves $r$ unchanged. This
invariance is exploited directly in Chapter~3, where it guarantees that
dividing every item by its scale maximum---a rescaling used to handle
mixed-format questionnaires---cannot change the correlation among items measured
on a common scale. Third, the \emph{sign} of $r$ encodes the direction of the
relationship, which is why reverse-keyed questionnaire items (worded so that
agreement indicates \emph{less} of the trait) correlate \emph{negatively} with
same-trait items keyed in the usual direction; handling that sign correctly is
essential when aggregating correlations across item pairs.

A distinction that is central to this thesis, but easy to overlook, is between
correlation computed \emph{across respondents} and correlation computed
\emph{within a respondent}. The familiar use of Equation~\ref{eq:pearson} treats
each respondent as one unit and each pair of items as the two variables $X$ and
$Y$; this \emph{sample-level} correlation describes how two items relate in the
population. One can instead fix a single respondent and treat their answers to a
set of item pairs as the data, computing a correlation that describes whether
\emph{that person} answered coherently across pairs that the sample says should
agree. This \emph{individual-level} or \emph{person-level} correlation is the
quantity ReReReRe estimates and tests; the psychometric-synonym index of
Chapter~1 is an early instance of the same idea
\parencite{meade2012identifying}.

% =====================================================================
\section{Likert Scales and the Measurement Model}
% =====================================================================

Most of the data this thesis concerns come from \emph{Likert items}
\parencite{likert1932technique}: statements to which respondents indicate
agreement on an ordered set of categories, typically five or seven points from
``strongly disagree'' to ``strongly agree.'' Strictly, a Likert response is
\emph{ordinal}---the categories are ordered, but the psychological distance
between ``agree'' and ``strongly agree'' need not equal that between ``neutral''
and ``agree.'' In practice, researchers routinely treat summed or averaged
Likert responses as \emph{interval} data and apply correlation- and
mean-based statistics to them, a simplification that is generally defensible for
scales with five or more categories that are reasonably symmetric.

The thesis follows this common practice, computing Pearson correlations on
Likert responses, with two consequences worth flagging now. First, because the
correlation in Equation~\ref{eq:pearson} is sensitive to the numerical spacing
of the response options, items administered on \emph{different} Likert ranges
(say, some on a 1--5 scale and others on a 1--7 scale) can contribute unequally
to an individual-level correlation; Chapter~3 addresses this with a rescaling
step. Second, a respondent who selects the same category for every item has zero
variance, which makes the denominator of Equation~\ref{eq:pearson} undefined for
that person---a degenerate case that careless ``straightlining'' produces and
that any individual-level method must handle explicitly.

% =====================================================================
\section{Latent Variables and Factor Analysis}
\label{sec:efa}
% =====================================================================

Why should answers to different items correlate at all? The dominant answer in
psychometrics is the \emph{common-factor model}, which traces to
\textcite{spearman1904general} and was systematized by
\textcite{thurstone1947multiple}. The model posits that the observed responses
are generated by a small number of unobserved \emph{latent factors}---traits
such as extraversion or conscientiousness---plus item-specific noise. Items that
load on the same factor correlate because they share that common cause; items on
different factors correlate weakly or not at all. Formally, the response of a
respondent to item $j$ is modelled as a weighted sum of the latent factors,
\begin{equation}
x_j \;=\; \lambda_{j1} f_1 + \lambda_{j2} f_2 + \cdots + \lambda_{jm} f_m
\;+\; \varepsilon_j ,
\label{eq:cfa}
\end{equation}
where $f_1,\dots,f_m$ are the latent factors, the \emph{loadings}
$\lambda_{jk}$ quantify how strongly item $j$ depends on factor $k$, and
$\varepsilon_j$ is the part of the item unique to itself (specific variance plus
measurement error). The squared loadings summed across factors give the item's
\emph{communality}, the proportion of its variance explained by the common
factors.

\emph{Exploratory factor analysis} (EFA) works backwards from the observed
correlation matrix to estimates of the loadings and the number of factors,
without assuming in advance which items load on which factor
\parencite{fabrigar1999evaluating, costello2005best}. It is ``exploratory'' in
contrast to confirmatory factor analysis, where the loading pattern is specified
ahead of time and tested. EFA matters for this thesis in two ways. The
\emph{generative} use comes first: the simulation studies of later chapters
create artificial questionnaires from a known factor model of the form in
Equation~\ref{eq:cfa}, so that the ground-truth structure is available for
evaluation. The \emph{analytic} use comes second: ReReReRe relies on the fact
that attentive respondents reproduce the sample's factor structure in their own
answers, and one variant of the method uses EFA to choose which item pairs to
score.

A practical question in any factor analysis is \emph{how many factors to
retain}. Retaining too few collapses distinct constructs; too many fits noise.
Two classical heuristics are \textcite{cattell1966scree}'s \emph{scree test},
which looks for an ``elbow'' in the plot of successive eigenvalues, and the
eigenvalue-greater-than-one rule. Both are now regarded as unreliable, and the
recommended procedure is \emph{parallel analysis} \parencite{horn1965rationale}:
one generates many random data sets with the same number of variables and
respondents but no real factor structure, computes their eigenvalues, and
retains only the factors whose eigenvalues in the real data exceed what random
data produce by chance. Parallel analysis is, in effect, a permutation-style
calibration of the factor count---a recurring logic in this thesis---and is the
factor-retention method ReReReRe uses internally when it needs an estimate of
dimensionality.

% =====================================================================
\section{Reliability and Internal Consistency}
% =====================================================================

A measurement is \emph{reliable} to the extent that it is reproducible---that
repeated measurement of the same unit yields similar values. For a multi-item
scale, the most widely reported reliability index is \emph{coefficient alpha}
\parencite{cronbach1951coefficient}, which summarizes the average inter-item
correlation, adjusted for the number of items, into a single number between
$0$ and $1$. High alpha indicates that the items ``hang together,'' covarying as
they should if they measure a common construct.

Reliability connects directly to careless responding. Because careless answers
are unrelated to the construct, they behave like added noise: they raise the
$\varepsilon_j$ term in Equation~\ref{eq:cfa}, lower communalities, weaken
inter-item correlations, and therefore \emph{depress} reliability
\parencite{huang2012detecting}. A drop in alpha when a scale is administered to
a careless-contaminated sample is one of the symptoms motivating detection in
the first place. Conversely, identifying and removing careless respondents
typically \emph{raises} the reliability and sharpens the factor structure of the
remaining data---the practical payoff that Chapter~5 evaluates.

% =====================================================================
\section{Permutation and Randomization Tests}
\label{sec:perm}
% =====================================================================

Classical hypothesis tests compare an observed statistic to a theoretical null
distribution---the $t$, $F$, or $\chi^2$ distribution---derived under
assumptions such as normality. When those assumptions are doubtful, a
\emph{permutation} (or \emph{randomization}) \emph{test} offers an alternative
that builds the null distribution \emph{from the data themselves}
\parencite{good2005permutation}. The logic is simple and assumption-light: if
the null hypothesis is true, then the observed arrangement of the data is just
one of many equally likely arrangements obtainable by shuffling. One therefore
repeatedly shuffles the data in a way consistent with the null, recomputes the
statistic on each shuffle, and reads off where the observed value falls in the
resulting distribution. A value far in the tail is unlikely under the null and
constitutes evidence against it.

This is precisely the inferential principle behind ReReReRe. The method asks, of
each respondent, whether their coherence across the item pairs that the sample
deems related is higher than would be expected if their answers carried no such
structure. The ``no structure'' null is realized by computing the same coherence
over \emph{randomly chosen} item pairs, repeated many times, to form a
per-respondent null distribution. The respondent's observed coherence is then
expressed as a \emph{z-score} relative to this distribution,
\begin{equation}
z \;=\; \frac{C_{\text{obs}} - \operatorname{mean}(C_{\text{rand}})}
{\operatorname{sd}(C_{\text{rand}})},
\label{eq:z}
\end{equation}
where $C_{\text{obs}}$ is coherence over the genuinely coupled pairs and
$C_{\text{rand}}$ is the coherence over random pairs. A high z means the
respondent is far more coherent than chance---attentive; a low or negative z
means they are no more coherent than random pairs would predict---potentially
careless. Building the baseline per respondent makes the method
\emph{self-calibrating}: it does not assume any particular distribution for the
coherence statistic, sidestepping exactly the normality problems that, as
Chapter~1 noted, make the chi-square cutoff for Mahalanobis distance unreliable
on Likert data.

% =====================================================================
\section{The Bootstrap}
\label{sec:boot}
% =====================================================================

A close cousin of the permutation test is the \emph{bootstrap}
\parencite{efron1979bootstrap, efron1993introduction}, a resampling method for
quantifying the uncertainty of an estimate. Where a permutation test shuffles to
simulate a null hypothesis, the bootstrap \emph{resamples with replacement} to
simulate the sampling variability of a statistic. Given a sample of $n$ units,
one draws many new samples of size $n$ by sampling the original units with
replacement (so some appear multiple times and others not at all), recomputes
the statistic of interest on each resample, and uses the spread of these
recomputed values as an estimate of the statistic's sampling distribution.

The bootstrap is valuable precisely when a statistic is too complicated for a
closed-form standard error. This thesis uses it for exactly such a case. To
quantify how much ReReReRe contributes to the ensemble, Chapter~4 reports the
\emph{difference} in Matthews correlation coefficient between the full ensemble
and an ablated ensemble without ReReReRe. No textbook formula gives the standard
error of a difference of two MCCs, so a \emph{percentile bootstrap} confidence
interval is used: respondents are resampled with replacement several hundred
times, the difference is recomputed on each resample, and the $2.5$th and
$97.5$th percentiles of the resulting distribution form a 95\% interval. If that
interval excludes zero, the contribution is statistically distinguishable from
no contribution---the form in which the thesis's central ablation claim is
stated.

% =====================================================================
\section{Supervised Classification and the Random Forest}
\label{sec:rf}
% =====================================================================

The ensemble at the heart of Chapters~4 and~5 is a \emph{supervised classifier}:
a model trained on examples whose true labels (careless or attentive) are known,
which then predicts labels for new respondents. The specific classifier is the
\emph{Random Forest} \parencite{breiman2001random}, built from \emph{decision
trees}. Because many psychologists know logistic regression but not tree
ensembles, this section develops the Random Forest from the ground up; a fuller
treatment appears in standard references \parencite{hastie2009elements,
james2013introduction}.

\subsection{Decision trees}

A \emph{decision tree} is a sequence of yes/no questions about the predictor
variables, arranged as a branching flowchart \parencite{breiman1984cart}. To
classify a respondent, one starts at the top and answers each question---``is
the ReReReRe z-score below $-0.5$?'', then ``is the longstring above $12$?''---
following the corresponding branch until reaching a terminal \emph{leaf}, which
carries a predicted class (or class probability). The tree is \emph{grown} from
training data by repeatedly choosing, at each node, the question that best
separates careless from attentive respondents according to a purity criterion
such as the Gini index. A single tree has an appealing property---its decisions
are transparent and can interact predictors automatically (the second question
asked depends on the answer to the first)---but also a serious flaw: it is
\emph{high-variance}. Small changes in the training data can change which
question is chosen near the top and thereby reshape the entire tree.

\subsection{From one tree to a forest}

The Random Forest tames this instability by averaging \emph{many}
de-correlated trees \parencite{breiman2001random}. Two sources of randomness
make the trees differ. First, each tree is grown on a \emph{bootstrap resample}
of the training respondents (the resampling idea of Section~\ref{sec:boot}),
so each tree sees a slightly different data set---a strategy called
\emph{bagging}, for bootstrap aggregating. Second, at each node only a
\emph{random subset} of the predictors is considered as candidate questions
(typically about $\sqrt{p}$ of the $p$ predictors), which prevents every tree
from being dominated by the same one or two strong predictors. The forest's
prediction is the average of the predictions of its trees---in our setting, the
proportion of trees that vote ``careless,'' interpreted as a predicted careless
probability. Averaging many noisy but roughly unbiased trees reduces variance
without inflating bias, which is why the forest predicts far more stably than any
single tree.

\subsection{Why a forest can beat logistic regression}

Logistic regression models the log-odds of the outcome as a \emph{linear,
additive} function of the predictors. That is an excellent model when the truth
is close to linear and additive, but careless detection is not such a case. The
predictors interact: a respondent with both a very low ReReReRe z-score
\emph{and} an unusually low intra-individual response variability is far more
suspicious than the additive combination of the two signals would imply, because
that conjunction is the signature of straightlining. A Random Forest captures
such interactions automatically through the nested structure of its splits,
without the analyst having to specify interaction terms in advance. In the
simulation studies of later chapters, the forest outperforms logistic regression
on every evaluation metric, with the largest margins on the hardest cases (short
questionnaires and low careless base rates).

A further benefit is \emph{variable importance}: by measuring how much
predictive accuracy degrades when a predictor is randomly permuted, a Random
Forest ranks the predictors by their contribution \parencite{strobl2008conditional}.
This provides a principled way to ask which detectors carry the signal---an
analysis the thesis uses to show that ReReReRe and a small set of auxiliary
indices account for most of the ensemble's power.

% =====================================================================
\section{Cross-Validation}
\label{sec:cv}
% =====================================================================

A classifier evaluated on the same data used to train it will look better than
it really is, because it can memorize idiosyncrasies of the training sample that
do not generalize---a problem called \emph{optimism} or \emph{overfitting}.
\emph{Cross-validation} estimates out-of-sample performance without requiring a
separate test set \parencite{stone1974cross, kohavi1995study}. In
\emph{$k$-fold} cross-validation, the data are partitioned into $k$ roughly
equal folds (here $k=5$). The model is trained on $k-1$ folds and used to
predict the held-out fold; this is repeated $k$ times so that each fold serves
once as the held-out set. Every respondent thus receives a prediction from a
model that did \emph{not} see them during training. These \emph{out-of-fold}
predictions are concatenated and used for all performance estimates, giving an
honest picture of how the classifier would behave on new respondents drawn from
the same population. The folds are \emph{stratified}---constructed to preserve
the careless base rate within each fold---so that no fold is unrepresentative.

% =====================================================================
\section{Evaluating a Binary Classifier}
\label{sec:metrics}
% =====================================================================

Once a classifier assigns each respondent a predicted careless probability,
turning those probabilities into a binary decision requires a \emph{threshold},
and judging the decisions requires \emph{metrics}. Both rest on the
\emph{confusion matrix} (Table~\ref{tab:confusion}), which cross-tabulates the
predicted label against the truth into four counts: true positives (TP, careless
correctly flagged), true negatives (TN), false positives (FP, attentive
respondents wrongly flagged), and false negatives (FN, careless respondents
missed).

\begin{table}[t]
\centering
\caption{The confusion matrix for careless detection. ``Positive'' denotes a
respondent flagged as careless.}
\label{tab:confusion}
\small
\begin{tabular}{l c c}
\toprule
 & \textbf{Truly careless} & \textbf{Truly attentive} \\
\midrule
\textbf{Flagged careless}  & True positive (TP)  & False positive (FP) \\
\textbf{Flagged attentive} & False negative (FN) & True negative (TN)  \\
\bottomrule
\end{tabular}
\end{table}

From these four counts a family of metrics is derived. \emph{Sensitivity}
(recall, true-positive rate) is the proportion of truly careless respondents
that are flagged, $\text{TP}/(\text{TP}+\text{FN})$; \emph{specificity} is the
proportion of attentive respondents correctly left unflagged,
$\text{TN}/(\text{TN}+\text{FP})$. \emph{Precision} (positive predictive value)
is the proportion of flagged respondents who are genuinely careless,
$\text{TP}/(\text{TP}+\text{FP})$. The \emph{F1 score} is the harmonic mean of
precision and recall. These are summarized in Table~\ref{tab:metrics}.

\begin{table}[t]
\centering
\caption{Principal evaluation metrics derived from the confusion matrix.}
\label{tab:metrics}
\small
\begin{tabular}{L{0.30\linewidth} L{0.40\linewidth} L{0.20\linewidth}}
\toprule
\textbf{Metric} & \textbf{Definition} & \textbf{Range} \\
\midrule
Sensitivity (recall) & $\text{TP}/(\text{TP}+\text{FN})$ & $0$ to $1$ \\
Specificity & $\text{TN}/(\text{TN}+\text{FP})$ & $0$ to $1$ \\
Precision (PPV) & $\text{TP}/(\text{TP}+\text{FP})$ & $0$ to $1$ \\
F1 & harmonic mean of precision and recall & $0$ to $1$ \\
MCC & Equation~\ref{eq:mcc} & $-1$ to $+1$ \\
\bottomrule
\end{tabular}
\end{table}

The thesis adopts the \emph{Matthews correlation coefficient} (MCC) as its
primary metric \parencite{matthews1975comparison, chicco2020advantages}. The MCC
is, in effect, the Pearson correlation between the predicted and the true binary
labels,
\begin{equation}
\text{MCC} \;=\;
\frac{\text{TP}\cdot\text{TN} - \text{FP}\cdot\text{FN}}
{\sqrt{(\text{TP}+\text{FP})(\text{TP}+\text{FN})
(\text{TN}+\text{FP})(\text{TN}+\text{FN})}} ,
\label{eq:mcc}
\end{equation}
ranging from $-1$ (perfect disagreement) through $0$ (chance) to $+1$ (perfect
agreement). Its decisive advantage over accuracy and F1 is that it uses all four
cells of the confusion matrix and is therefore robust to \emph{class
imbalance}, which careless detection always involves: when only, say, 10\% of
respondents are careless, a useless classifier that flags no one achieves 90\%
accuracy, and even F1 can be inflated, whereas the MCC of such a classifier is
$0$ \parencite{chicco2020advantages}. A high MCC requires the classifier to do
well on \emph{both} classes simultaneously.

Two \emph{threshold-free} metrics complete the picture by summarizing
performance across all possible thresholds at once. The \emph{area under the
receiver operating characteristic curve} (AUC) is the probability that a randomly
chosen careless respondent receives a higher predicted probability than a
randomly chosen attentive one \parencite{hanley1982meaning}; $0.5$ is chance and
$1$ is perfect. The \emph{area under the precision--recall curve} (AUPRC) is a
companion measure that, by ignoring the typically very large true-negative cell,
is more informative than AUC under heavy class imbalance
\parencite{saito2015precision}. Because AUC and AUPRC do not depend on a chosen
cutoff, they isolate the quality of the underlying ranking from the separate
question of where to set the threshold.

That separate question---\emph{which} threshold to use---is itself consequential,
as Chapter~1 stressed. A threshold defines an \emph{operating point}, a specific
trade-off between sensitivity and specificity; moving it trades missed careless
respondents against falsely flagged attentive ones. Chapter~5 returns to the
problem of choosing an operating point without access to ground-truth labels,
which is the form the threshold problem takes in real deployment.

% =====================================================================
\section{Summary}
% =====================================================================

The pieces assembled here recur throughout the thesis as a connected toolkit
rather than a list. Correlation, computed at the level of the individual rather
than the sample, is what ReReReRe measures; the common-factor model explains why
that correlation exists for attentive respondents and parallel analysis
estimates its dimensionality; the permutation test turns the measurement into a
self-calibrating per-respondent z-score; the bootstrap puts a confidence
interval on the detector's contribution; the Random Forest combines that
detector with auxiliary indices, cross-validation estimates the combination's
performance without optimism, and the Matthews correlation coefficient---together
with AUC and AUPRC---reports how good that performance is under the class
imbalance intrinsic to careless detection. With these concepts in place,
Chapter~3 can introduce the ReReReRe detector itself.

\printbibliography

\end{document}
